% Chapter 5 -- Rubric: PO4 (4.1)
\chapter{Investigation and Evaluation}
\label{ch:evaluation}

\section{Experimental Investigation}
\label{sec:investigation}

This is a development-focused project, so the investigations reported here are
those that validated --- or refuted --- specific design decisions. Four are
presented. Each states the question it was asked to settle, the method, and the
conclusion actually drawn, including the two cases where the investigation
overturned the team's initial position.

\subsection{Investigation 1: is local model inference viable for content generation?}

\textbf{Question.} Can the question bank be generated on the team's own hardware
using locally hosted open-weight models, avoiding any external API dependency
and any cost?

\textbf{Method.} A generation pipeline was implemented against locally served
quantised models (7B and 3B parameter classes) on CPU-only laptops, and measured
on real specifications drawn from the ontology, with prompt sizes of roughly
4--6\,k tokens and completions of roughly 3\,k tokens per call.

\textbf{Findings.} Two independent failures. On throughput, the measured
per-call latency extrapolated to multiple weeks of continuous computation for a
single full pass over the 2{,}580-specification space then in scope --- longer
than the remaining project schedule, on machines also needed for development. On
correctness, a subtler defect appeared: models emit \LaTeX{} commands such as
\verb|\tan| and \verb|\frac| with single backslashes, and a strict JSON parser
does not reject these --- for the subset of letters that happen to be legal JSON
escapes it silently consumes the backslash and substitutes a control character.
The corruption therefore produced no error at all, and was found only by reading
output.

\textbf{Conclusion.} The local approach was abandoned on throughput grounds, and
a hosted API adopted with key pooling to raise the rate ceiling. The escape
defect was not specific to local models --- it recurred with the hosted model at a
measured 13.4\% of one early production batch --- so an escape-repair pass was
made unconditional in both the ingest and generation paths. This is the clearest
case in the project of an investigation changing the architecture: the design in
\Cref{sec:generation-pipeline} exists because this one failed.

\subsection{Investigation 2: is the ontology's granularity sufficient to support the gating policy?}

\textbf{Question.} The conjunctive gate in \Cref{eq:gating} acts on a skill's
direct prerequisites. Does the first ontology supply enough structure for that
gate to do anything?

\textbf{Method.} Structural analysis of the 430-skill graph: skills per topic,
syllabus coverage, and the fraction of skills with no prerequisite edge at all.

\textbf{Findings.} Two deficits, both fatal to the gate. The graph averaged
fewer than four skills per topic, and a topic holding three skills cannot
localise a failure \emph{within} itself --- which is the system's entire
purpose. And a substantial fraction of skills had no incoming edge at all; a
gate with no parents to test degenerates to the ungated case, so for those
skills \Cref{eq:gating} was doing nothing.

Note that the binding measure here is coverage and connectedness, not the raw
edges-per-skill ratio. That ratio is similar before and after the rebuild
(1.1 against 1.03), and deliberately so: the rebuild added roughly four times
as many skills, so holding the ratio while eliminating isolated skills entirely
means the edges are distributed across the graph rather than concentrated in a
few dense topics. A higher ratio achieved by piling edges onto already-connected
skills would not have fixed either deficit above.

\textbf{Conclusion.} The ontology, not the engine, was the binding constraint.
This motivated the full-corpus rebuild (\Cref{sec:ontology-rebuild}), which
raised the graph to 1{,}688 skills and 1{,}743 edges with every one of the 183
syllabus topics populated and no isolated skills. \Cref{tab:ontology-effects}
reports the rebuild's structural effect.

\begin{table}[htbp]
  \centering
  \caption{Structural effect of the ontology rebuild and its editorial pass}
  \label{tab:ontology-effects}
  \begin{tabular}{lrr}
    \toprule
    Property & Before editorial pass & After \\
    \midrule
    Syllabus topics holding no skill & 54 & 0 \\
    Unresolved prerequisite references & 46 & 0 \\
    Isolated skills (no edge at all) & 6 & 0 \\
    Topics with no internal learning order & 5 & 0 \\
    Near-orphan fragments (2--6 skills) & 91 & 32 \\
    Connected components & 187 & 87 \\
    Largest connected component & 203 & 482 \\
    Maximum DAG depth & 7 & 9 \\
    \bottomrule
  \end{tabular}
\end{table}

\subsection{Investigation 3: can prerequisite edges be proposed mechanically?}

\textbf{Question.} After the rebuild, 91 fragments of 2--6 skills remained
disconnected from the main graph. Can the missing edges be proposed
automatically rather than authored by hand?

\textbf{Method.} Two scoring functions were implemented and their proposals
read individually against the subject matter. The first scored candidate parents
by token overlap with the orphaned skill. The second scored by
\emph{topic hub-ness} --- preferring the candidate at the lowest Bloom tier with
the most dependents.

\textbf{Findings.} Token overlap produced roughly 40\% usable proposals. Its
failures were systematic rather than random: it proposed lateral links between
peers that merely share vocabulary --- zero-order kinetics ``requiring''
first-order kinetics, the cosine rule ``requiring'' the sine rule --- and in one
case proposed a pair in both directions, which would have closed a cycle.
Hub-based scoring raised usable proposals to roughly two thirds. Of 83 proposals
finally reviewed, 48 were applied (several with the direction reversed from what
was proposed) and 35 were rejected outright, their fragments left deliberately
disconnected rather than wrongly connected.

\textbf{Conclusion.} Mechanical proposal is a useful accelerator and an
unacceptable authority. A wrong prerequisite edge is worse than a missing one,
because \Cref{eq:pullup} propagates along it and will credit a learner for a
skill they have never demonstrated. The 32 fragments that survive are recorded
as requiring a domain expert rather than another heuristic.

\subsection{Investigation 4: do automated quality heuristics over-flag?}

\textbf{Question.} Two validators were written to find defects in the generated
ontology: one flagging skills that fuse two actions, and one flagging ``Bloom
inversions'' --- edges whose prerequisite sits at a \emph{higher} Bloom tier than
the skill requiring it. Are their outputs trustworthy?

\textbf{Method.} Every flag from both checks was read against the underlying
skill or edge.

\textbf{Findings.} The multi-action check flagged 173 skills, of which over half
were false positives, and was replaced by a stricter criterion yielding 47
candidates --- of which 20 were genuine splits and 27 were rejected on reading.
The Bloom-inversion check flagged 147 edges; hand review found \emph{3}
genuinely reversed. The check's premise was wrong, not the data: Bloom tier
measures cognitive demand, not teaching order, and a low-tier fact \emph{about}
an operation is routinely learned after the operation itself. For example,
recalling that a determinant with two equal rows is zero (Remember) legitimately
depends on being able to evaluate a determinant (Apply).

\textbf{Conclusion.} The inversion check was downgraded from error to advisory
rather than deleted, with its reasoning recorded so the next contributor does not
reinstate it. The transferable finding is that on this project mechanical
heuristics over-flagged by roughly the same large factor twice, and that the
correct response is to calibrate a heuristic against hand review before treating
its output as a defect list.

\section{Performance Evaluation}
\label{sec:evaluation}

\subsection{Setup and acceptance criteria}

The engine is evaluated by an automated end-to-end harness that runs the real
diagnostic, mastery-update, topic-practice and spillover code paths against an
in-memory stand-in for the database, seeded with a synthetic learner and a
question bank drawn from the compiled ontology. It requires no credentials and
no network, which is what makes it runnable by any developer before a merge.
Acceptance is binary per check: all checks must pass.

\subsection{Results}

The harness reports \textbf{18 checks passed, 0 failed}, reproduced on the
current tree against the adopted 1{,}688-skill ontology. \Cref{tab:smoke} lists
what each check establishes.

One property of the harness should be read before its results are interpreted.
It selects a section automatically and sizes the diagnostic to that section's
skill count, so the absolute figures it prints depend on which section it picks
--- on the current tree that is a small one (eight skills). The checks therefore
establish that each mechanism \emph{runs and holds its invariants}; they are not
a measurement of how far evidence propagates in a large section, which would
require a section with real depth and a populated question bank. The propagation
guarantee is carried by the invariant check, not by counting how many skills
moved.

\begin{table}[htbp]
  \centering
  \caption{End-to-end verification harness: checks and what each establishes}
  \label{tab:smoke}
  \begin{tabularx}{\textwidth}{>{\hsize=0.42\hsize}Y>{\hsize=0.82\hsize}Y>{\hsize=1.76\hsize}Y}
    \toprule
    Area & Check & Establishes \\
    \midrule
    Ontology & Graph builds (1{,}688 skills); catalogue topics all resolve (183); no legacy topic codes remain & The served ontology is internally consistent and the catalogue cannot reference a topic that does not exist \\
    Question search & Tier-1 filter matches by topic \emph{code}; filtering by display label matches nothing & Regression guard for the defect in \Cref{sec:revisions} \\
    Diagnostic & Session starts, sized to the section (8 skills in the section exercised); question carries its subject; topic shown as label (\texttt{PHY1\_ERRORS} $\rightarrow$ ``Errors and Significant Figures''); every question answered & The adaptive loop runs to completion, sizes itself to the section, and returns learner-presentable metadata (FR3, FR12) \\
    Mastery & Mastery unlocked after diagnostic; written to \texttt{user\_skill}; section state reports subject & The update path persists per-skill mastery and releases the diagnostic gate (FR4, FR7) \\
    Invariant & No ancestor sits below its descendant & \Cref{eq:pullup} holds globally after a full session --- the correctness property in NFR1 \\
    Mastery view & View unlocked; graph returns nodes and edges for the section; table rows carry subject & The visualisation contract holds (FR10) \\
    Practice & Topic practice serves and accepts questions & Post-diagnostic flow works (FR8) \\
    Spillover & Fires on repeated wrong answers, reporting its triggering rule (rule A, repeated missing prerequisite) & The prerequisite-detection policy triggers on the intended evidence (FR9) \\
    \bottomrule
  \end{tabularx}
\end{table}

The ontology is separately validated by a structural checker, which reports
\textbf{zero errors}: the graph is acyclic, every one of the 183 topics holds at
least one skill, no skill is isolated, and no prerequisite reference is
unresolved.

\subsection{Evaluation against requirements}

\Cref{tab:traceability} traces every functional and non-functional requirement
from \Cref{ch:requirements} to the evidence that it is met, or to the deviation
recorded in \Cref{sec:revisions} where it is not. Two entries are not claims of
success: FR13 is partially met and NFR2 was never measured.

\begin{table}[htbp]
  \centering
  \caption{Requirements traceability and status}
  \label{tab:traceability}
  \begin{tabularx}{\textwidth}{>{\hsize=0.3\hsize}Y>{\hsize=0.45\hsize}Y>{\hsize=2.25\hsize}Y}
    \toprule
    Req. & Status & Evidence or deviation \\
    \midrule
    FR1 & Met, reduced & Account creation and sign-in work, but without a password (\Cref{sec:deviations-auth}) \\
    FR2 & Met & Catalogue endpoint serves ontology-derived courses, sections and topics; verified by harness \\
    FR3 & Met & Diagnostic runs to a full 30-question session in the harness \\
    FR4 & Met & Bloom-conditioned update implemented per \Cref{tab:bloom-params}; exercised on every answered question \\
    FR5 & Met & Ancestor pull-up implemented; the invariant that no ancestor sits below its descendant is verified explicitly after a full session \\
    FR6 & Met & Conjunctive gate implemented per \Cref{eq:gating} \\
    FR7 & Met & Mastery view locked before, unlocked after the diagnostic \\
    FR8 & Met & Topic practice served and answered questions in the harness \\
    FR9 & Met & Spillover fired and reported its triggering rule \\
    FR10 & Met & Mastery graph (nodes and edges for the section) and table both returned \\
    FR11 & Met & Retake endpoint resets section mastery and clears sessions \\
    FR12 & Met & Mathematics renders through KaTeX; literal newline handling fixed during implementation \\
    FR13 & \textbf{Partially met} & Pipeline is complete and validated, but the bank covers only part of the specification space (\Cref{sec:eval-limits}) \\
    FR14 & Met & Interface strings available in English and Bangla \\
    \midrule
    NFR1 & Met & Acyclicity enforced at insertion; ancestor invariant checked by harness; validator reports zero errors \\
    NFR2 & \textbf{Not measured} & No latency benchmark was run (\Cref{sec:eval-limits}) \\
    NFR3 & Met & Fallback ladder exercised; a missing item marks the skill tested and advances \\
    NFR4 & Met & Ontology compiled from source by script; generated files are not hand-edited \\
    NFR5 & Met & Harness runs with no credentials and no network \\
    NFR6 & Met & Entire system developed and run on CPU-only laptops \\
    NFR7 & Met & Bangla content with \LaTeX{} mathematics end to end \\
    NFR8 & Met & Credentials held in a git-ignored environment file, server-side only \\
    \bottomrule
  \end{tabularx}
\end{table}

\subsection{Evaluation against objectives}

O1--O5, O7 and O8 are met as evidenced above.

\textbf{O6 (question bank) is partially met}, and the shortfall is worth
quantifying rather than leaving as an adjective. The pipeline is built,
validated, and has generated and loaded real verified items; what is incomplete
is coverage of the specification space. Two figures bound the gap. Against the
legacy 430-skill ontology the space was $430 \times 6 = 2{,}580$
(skill, Bloom level) specifications, and generation had reached roughly a fifth
of it when the rebuild superseded that target. Against the now-adopted
1{,}688-skill ontology the space is $1{,}688 \times 6 = 10{,}128$
specifications, and because the existing items are keyed to legacy skill
identifiers they do not transfer: coverage against the served ontology therefore
starts close to zero and the generation run is outstanding in full. This is the
same debt identified in \Cref{sec:deviations-adoption}, and it is the single
largest piece of remaining work (\Cref{sec:future}).

\textbf{O3 is met in implementation but unvalidated in effect} --- the
propagation is provably consistent, but whether the resulting estimates
correspond to real learner knowledge is exactly the question the next subsection
says was not answered.

\subsection{Limitations of this evaluation}
\label{sec:eval-limits}

The evaluation above is structural and functional. It establishes that the
system computes what it was designed to compute and that its invariants hold. It
does \textbf{not} establish that the system helps anyone learn, and the report
makes no such claim. Specifically:

\begin{itemize}
  \item \textbf{No human participants.} There was no user study, no classroom
    trial, and no usability testing with members of the target cohort. Every
    claim about learner benefit in this report is an argument from design, not a
    measurement.
  \item \textbf{No fitted parameters.} $P(L_0)$, $P(T)$ and the Bloom-conditioned
    $P(G), P(S)$ values are pedagogical priors, not estimates. The literature is
    explicit that fitted and individualised parameters outperform fixed ones
    \cite{yudelson2013individualized,baker2008more}; this system cannot do that
    until it has response data, which it cannot have before deployment.
  \item \textbf{No predictive validation.} The standard quantitative check on a
    knowledge-tracing model --- does it predict the next response better than a
    baseline, measured by AUC on held-out data --- was not performed, because
    there is no held-out learner data.
  \item \textbf{No performance benchmark.} NFR2 was not measured; the system is
    subjectively responsive in use, which is not evidence.
  \item \textbf{Ontology unreviewed by domain experts.} The rebuilt graph is
    structurally valid but its subject-matter correctness has been verified only
    by the team. 32 disconnected fragments and a body of hand-authored
    prerequisite edges are explicitly awaiting review by a chemist and a
    physicist.
  \item \textbf{Question bank coverage is partial}, so a full-coverage session is
    not yet reachable for every topic. Two consequences follow that the harness
    cannot surface, because it seeds synthetic items rather than reading the real
    bank. First, the Bloom ladder degrades: the engine requests a question one
    band above the learner's current band, and where the bank holds no item at
    that level the nearby-Bloom fallback
    ($[0,+1,-1,+2,-2,+3,-3]$) supplies whatever exists, so the cognitive
    progression the design intends is weaker in practice than in specification.
    Second, the graph-nearby search tiers are exercised far less than the
    fallback, for the same reason. Both resolve with coverage, not with code.
  \item \textbf{The harness measures mechanism, not magnitude.} It sizes each
    run to the section it selects, so its printed figures vary with that choice
    and are evidence that each path executes and its invariants hold --- not a
    measurement of how far evidence propagates across a large, well-connected
    section.
\end{itemize}

\section{Deviations and Design Revisions}
\label{sec:revisions}

\subsection{Topic filter matched no rows (found, analysed, fixed)}

\textbf{Deviation.} The diagnostic's first two search tiers returned nothing,
ever. Every question was silently arriving from the untargeted fallback, so
topic-targeted and graph-nearby selection --- two thirds of the selection design
--- were dead code in practice.

\textbf{Cause.} The \texttt{questions.topic} column is a foreign key to
\texttt{ontology\_topics} and therefore holds a topic \emph{code}
(\texttt{MAT\_MATRIX}), while the search filtered it by the catalogue
\emph{display label} (``Matrices and Determinants''). The two can never be
equal, so the filter matched no row. Nothing failed loudly: a query returning
zero rows is a valid query, and the fallback dutifully supplied a question, so
sessions looked healthy.

\textbf{Revision.} The filter was corrected to compare codes, and --- more
importantly --- two checks were added to the harness: one asserting that a
code-based filter matches, and one asserting that a label-based filter matches
\emph{nothing}. The second check is the one that matters; it fails if anyone
reintroduces the confusion.

\subsection{Session state is not persistent}
\label{sec:deviations-session}

\textbf{Deviation.} Diagnostic and practice sessions live in process memory, so
a backend restart discards every in-flight session (NFR-adjacent, and a
production defect).

\textbf{Analysis.} The session objects hold graph traversal state that was not
designed for serialisation. Persisting them properly means either a session
store or reconstructing session state from the durable mastery table on demand.

\textbf{Decision.} Not revised, deliberately. The deployment target for this
project is a single demonstration process where the failure mode is a lost
session rather than lost mastery --- mastery itself \emph{is} persisted. The
correct fix is recorded as production hardening rather than disguised.

\subsection{Authentication is username-only}
\label{sec:deviations-auth}

\textbf{Deviation.} A learner signs in with a username and no password; the
client holds the resulting identity in browser storage. This does not meet the
access-control expectations of \Cref{sec:standards}.

\textbf{Analysis.} The system stores no data that is sensitive in itself beyond
performance history, and the prototype's users are the team. Authentication was
therefore deprioritised against the ontology work, which was on the critical
path for every other subsystem.

\textbf{Decision.} Not revised within the project. It is the first item of
production hardening (\Cref{sec:future}), and is stated here rather than
presented as a design choice, because it is not one --- it is an accepted debt.

\subsection{Ontology adoption, and the coverage debt it exposes}
\label{sec:deviations-adoption}

\textbf{Deviation.} For most of the project the application served the 430-skill
legacy ontology while the validated 1{,}688-skill rebuild sat beside it,
unadopted. Adoption has since been carried out: the compiled catalogue in
\texttt{Backend/tree\_data/} is now generated from the rebuild --- 1{,}688
skills across all 183 syllabus topics --- and the generation pipeline's ontology
source has been repointed to the rebuild so that the two cannot drift apart.

\textbf{Analysis.} Adoption was deliberately held back rather than forgotten,
because it is not merely a file swap. The existing question bank is keyed to
legacy skill identifiers, so switching the catalogue does not carry the content
across with it: a system that knows about 1{,}688 skills but holds questions for
a fraction of them can diagnose more finely in principle while being able to ask
about less in practice. Adopting therefore converts a \emph{structural} deficit
into a \emph{coverage} deficit --- the right trade, because coverage is bought
by a generation run whereas structure had to be rebuilt by hand, but it is a
debt and is recorded as one.

\textbf{Decision.} Adopt, then generate --- the order finally taken, with the
consequence stated explicitly rather than presented as complete. The outstanding
generation run against the enlarged specification space is quantified in
\Cref{sec:budget} and carried as the first item of
\Cref{sec:future}. Until it is complete, question-bank coverage against the
served ontology is the system's binding limitation, and \Cref{sec:eval-limits}
records it as such.
